% !TeX root = ../../Book5.tex
% 纲要标题：### 3.4 函子关系
% 纲要位置：工作记录/计划纲要.md，第 3809 行。
% 纲要细目（写作范围）：
% * 诱导的传递性
% * 张量恒等式
% * 有限指数情形诱导与余诱导的关系、右侧互反律的具体映射；无限指数时的陪集有限支集
% ---


\section{函子关系}
\label{b5:sec:0304}

诱导常常要和另一项构造接连使用：沿子群链诱导两次，把一个表示先张量再诱导，或研究映入诱导空间的同态。它们的相容性要由模同构表达，构造这些同构的办法是追踪张量或函数在连续操作中的去向。下面给出的映射与逆映射均在任意域上成立。

\subsection{诱导的传递性}
\label{b5:subsec:0304:01}

\begin{theorem}[诱导的传递性]\label{b5:thm:induction-transitivity}
	设 \(k\) 为域，\(K\le H\le G\) 为有限群的子群链，\(U\) 为左 \(k[K]\)-模。存在自然的左 \(k[G]\)-模同构
	\[
		\operatorname{Ind}_H^G\left(\operatorname{Ind}_K^H U\right) \cong \operatorname{Ind}_K^G U, \qquad a\otimes(b\otimes u)\longmapsto ab\otimes u \quad(a\in k[G],\ b\in k[H],\ u\in U).
	\]
\end{theorem}
\begin{proof}
	映射对 \(a,b,u\) 多线性。内层的关系 \(bc\otimes u=b\otimes cu\)，\(c\in k[K]\)，在目标中成为 \(abc\otimes u=ab\otimes cu\)。外层的关系 \(ah\otimes(b\otimes u)=a\otimes(hb\otimes u)\)，\(h\in k[H]\)，在目标中两边都等于 \(ahb\otimes u\)。所以映射良定义，且保持两端的 \(G\) 左乘作用。

	逆映射为 \(c\otimes u\mapsto c\otimes(1\otimes u)\)。对 \(d\in k[K]\)，有
	\[
		cd\otimes(1\otimes u) =c\otimes(d\otimes u) =c\otimes(1\otimes du),
	\]
	故它也保持张量关系。一个方向的复合显然为恒等；另一个方向使用 \(ab\otimes(1\otimes u)=a\otimes(b\otimes u)\)，同样恢复原元。对 \(K\)-同态 \(U\rightarrow U'\)，两条映射仅把最后一个分量送到其像，因此同构自然。
\end{proof}

选择 \(G/H\) 的左陪集代表 \(T\) 及 \(H/K\) 的左陪集代表 \(S\)，则 \(\{ts:t\in T,s\in S\}\) 是 \(G/K\) 的左陪集代表集。上述同构把 \(t\otimes(s\otimes u)\) 送到 \(ts\otimes u\)。具体地，若 \(gt=t'h\)、\(hs=s'k_0\)，其中 \(t'\in T\)、\(s'\in S\)、\(h\in H\)、\(k_0\in K\)，则两层诱导中的作用为
\[
	g\bigl(t\otimes(s\otimes u)\bigr) =t'\otimes\bigl(s'\otimes k_0u\bigr).
\]
而 \(gts=t's'k_0\)，故直接从 \(K\) 诱导时得到 \(t's'\otimes k_0u\)。两种坐标不仅有同样多的分量，而且把同一个 \(K\)-元素作用在 \(u\) 上，这才验证了群作用的相容性。

\subsection{张量恒等式}
\label{b5:subsec:0304:02}

若 \(V_1,V_2\) 为同一群的群代数模，域上张量积使用对角作用 \(g(v_1\otimes v_2)=gv_1\otimes gv_2\)。下面需要把这个作用与诱导张量积中的平衡关系分开。先看 \(H=\{1\}\)、\(W=k\) 的情形：同一个向量空间 \(k[G]\otimes_k V\) 在诱导表示中仅对群代数因子作用，公式是 \(x(g\otimes v)=xg\otimes v\)，而在两个表示的张量积上取对角作用，公式则是 \(x(g\otimes v)=xg\otimes xv\)。恒等映射一般不保持这两种作用；把 \(g\otimes v\) 送到 \(g\otimes gv\)，则 \(xg\otimes v\) 的像正好是 \(xg\otimes xgv\)。因此下面公式中的 \(gv\) 是为使两端作用相容而引入的，逆映射相应要用 \(g^{-1}v\)。

\begin{theorem}[张量恒等式]\label{b5:thm:tensor-identity}
	设 \(k\) 为域，\(G\) 为有限群，\(H\le G\)，\(W\) 为左 \(k[H]\)-模，\(V\) 为左 \(k[G]\)-模。以下同构对两个模自然：
	\begin{equation}\label{b5:eq:tensor-identity}
		\operatorname{Ind}_H^G\left(W\otimes_k\operatorname{Res}_H^G V\right) \cong (\operatorname{Ind}_H^G W)\otimes_k V.
	\end{equation}
	它在 \(g\in G\) 的纯张量上由
	\[
		g\otimes(w\otimes v)\longmapsto(g\otimes w)\otimes gv
	\]
	给出；右边使用 \(G\) 的对角作用。
\end{theorem}
\begin{proof}
	群元素是 \(k[G]\) 的基，所以给定公式先按 \(g\) 线性延伸，具体为
	\[
		\left(\sum_g c_gg\right)\otimes(w\otimes v) \longmapsto\sum_g c_g(g\otimes w)\otimes gv.
	\]
	不能把一般 \(a\in k[G]\) 直接放进右边的两个位置写成 \((a\otimes w)\otimes av\)：那样会产生两个求和指标和系数的乘积，不再是关于 \(a\) 的线性延伸。这里逐个群元素同时作用于两个因子，再对这些基元素的像求和。对 \(h\in H\)，映射在 \(gh\otimes(w\otimes v)\) 处的值为
	\[
		(gh\otimes w)\otimes ghv=(g\otimes hw)\otimes g(hv),
	\]
	这正是在 \(g\otimes(hw\otimes hv)\) 处的值，故保持 \(k[H]\) 上的平衡关系。对 \(x\in G\)，将 \(g\) 换成 \(xg\) 后的像为 \((xg\otimes w)\otimes xgv\)，等于原像受到 \(x\) 的对角作用，故它是 \(G\)-同态。

	逆映射在群元素上为
	\[
		(g\otimes w)\otimes v\longmapsto g\otimes(w\otimes g^{-1}v).
	\]
	其良定义需要检查 \(gh\otimes w=g\otimes hw\)。从前者出发，像为
	\[
		gh\otimes(w\otimes h^{-1}g^{-1}v) =g\otimes(hw\otimes g^{-1}v),
	\]
	与从后者出发相同。两种复合分别消去 \(g^{-1}g\) 与 \(gg^{-1}\)，故为恒等。\(H\)-同态与 \(G\)-同态分别保持相应的群作用，所以逐个代入这两条公式也验证了两变量自然性。
\end{proof}

这里的名称为\term[zhanglianghengdengshi]{张量恒等式}[tensor identity]，也常按其用途称投影公式。复数域的特征标取值满足张量积对应逐点相乘，故\eqref{b5:eq:tensor-identity}给出
\[
	\operatorname{Ind}_H^G\left(\psi\cdot\operatorname{Res}_H^G\chi\right) =(\operatorname{Ind}_H^G\psi)\cdot\chi.
\]
这个等式对任意复类函数 \(\psi,\chi\) 也成立：在\eqref{b5:eq:induced-character-all}中，\(\chi(x^{-1}gx)=\chi(g)\)，可直接将该因子提出求和。

\ProseParagraphBreak
例如令 \(H=\{1\}\)、\(W=k\)，可得 \(k[G]\otimes_k V\) 上的对角作用同构于 \(\dim_k V\) 个正则表示的直和（\(V\) 有限维时）。这不是半单性结论；\cref{b5:thm:tensor-identity}中的显式映射说明，即使 \(\operatorname{char}k\) 整除 \(|G|\)，该同构仍然成立。

\subsection{有限群情形诱导与余诱导的关系}
\label{b5:subsec:0304:03}

Frobenius 互反律将 \(\operatorname{Hom}_{k[G]}(\operatorname{Ind}_H^G W,V)\) 识别为 \(\operatorname{Hom}_{k[H]}(W,\operatorname{Res}_H^G V)\)。若给定的却是 \(\operatorname{Res}_H^G V\rightarrow W\)，就要寻找一个 \(G\)-模，使映入它的同态对应于这些 \(H\)-同态。下面的余诱导满足
\[
	\operatorname{Hom}_{k[G]}(V,\operatorname{Coind}_H^G W) \cong\operatorname{Hom}_{k[H]}(\operatorname{Res}_H^G V,W).
\]
这里也允许 \(G\) 为一般群：限制仍是只保留 \(H\) 的作用，诱导仍定义为 \(k[G]\otimes_{k[H]}W\)。群代数中的元素始终是群元素的有限线性组合，因此这两个定义不需要 \(G\) 本身有限；下面比较诱导与余诱导时才需要 \([G:H]<\infty\)。

\begin{definition}\label{b5:def:coinduction}
	设 \(k\) 为域，\(G\) 为群，\(H\le G\)，\(W\) 为左 \(k[H]\)-模。把 \(k[G]\) 以左乘看成左 \(k[H]\)-模，定义
	\[
		\operatorname{Coind}_H^G W \coloneqq\operatorname{Hom}_{k[H]}\bigl(k[G],W\bigr), \qquad (gf)(a)=f(ag).
	\]
	这里 \(g\in G\)、\(f\in\operatorname{Hom}_{k[H]}(k[G],W)\)、\(a\in k[G]\)。所得左 \(k[G]\)-模称为 \(W\) 的余诱导模。若 \([G:H]<\infty\) 且 \(W\) 有限维，则此模也是有限维的，称为 \(W\) 的\term[yuyoudaobiaoshi]{余诱导表示}[coinduced representation]。
\end{definition}
\symidx{\(\operatorname{Coind}_H^G\)：从 \(H\) 到 \(G\) 的余诱导}

该作用保持左 \(H\)-线性，因为 \(f(hag)=hf(ag)\)；它满足 \(\bigl(g_1(g_2f)\bigr)(a)=f(ag_1g_2)=\bigl((g_1g_2)f\bigr)(a)\)。等价地，余诱导空间由函数 \(f:G\rightarrow W\) 组成，满足
\begin{equation}\label{b5:eq:coinduced-functions}
	f(hx)=hf(x),\qquad (gf)(x)=f(xg).
\end{equation}
函数在群元素上的取值唯一线性延拓到群代数。这时使用的是右陪集 \(H\backslash G\)；它与前面的左陪集模型不可混用。右移自变量保留了 \(hx\) 左端的 \(h\)，所以 \(f(hxg)=hf(xg)\) 仍满足所需条件。若直接改成 \(f(gx)\)，在 \(hx\) 处就变成 \(f(ghx)\)，一般不能把 \(h\) 提到函数外；而且两次这样的左移会把 \(g_1,g_2\) 的乘法次序倒过来。

\begin{proposition}\label{b5:prop:coinduction-adjunction}
	设 \(k\) 为域，\(G\) 为群，\(H\le G\)，\(V\) 为左 \(k[G]\)-模，\(W\) 为左 \(k[H]\)-模。存在自然同构
	\[
		\operatorname{Hom}_{k[G]}(V,\operatorname{Coind}_H^G W) \cong\operatorname{Hom}_{k[H]}(\operatorname{Res}_H^G V,W).
	\]
	其正向映射为 \(F\mapsto\bigl(v\mapsto F(v)(1)\bigr)\)，逆映射把 \(u\) 送到 \(F_u(v)(x)=u(xv)\)。
\end{proposition}
\begin{proof}
	若 \(F\) 为 \(G\)-同态，则 \(F(hv)(1)=\bigl(hF(v)\bigr)(1)=F(v)(h)=hF(v)(1)\)，故取值于 \(1\) 得到 \(H\)-同态。反之，\(F_u(v)(hx)=u(hxv)=hu(xv)\)，所以 \(F_u(v)\) 满足\eqref{b5:eq:coinduced-functions}；并且 \(F_u(gv)(x)=u(xgv)=\bigl(gF_u(v)\bigr)(x)\)，故 \(F_u\) 为 \(G\)-同态。

	在 \(1\) 处求值恢复 \(u(v)\)；反向恢复 \(F\) 则使用 \(F(xv)(1)=\bigl(xF(v)\bigr)(1)=F(v)(x)\)。前复合 \(G\)-同态与后复合 \(H\)-同态均可直接移入取值公式，因此该对应对两个变量自然。
\end{proof}

\begin{theorem}\label{b5:thm:induction-coinduction}
	设 \(k\) 为任意域，\(G\) 为群，\(H\le G\) 且 \([G:H]<\infty\)。对每个左 \(k[H]\)-模 \(W\)，存在自然的 \(k[G]\)-模同构
	\[
		\operatorname{Ind}_H^G W\cong\operatorname{Coind}_H^G W.
	\]
	因而对每个左 \(k[G]\)-模 \(V\)，有自然同构
	\begin{equation}\label{b5:eq:frobenius-right-hom}
		\operatorname{Hom}_{k[G]}(V,\operatorname{Ind}_H^G W) \cong\operatorname{Hom}_{k[H]}(\operatorname{Res}_H^G V,W).
	\end{equation}
	取 \(G/H\) 的左陪集代表集 \(T\)，右边的同态 \(u\) 对应于
	\[
		F_u(v)=\sum_{t\in T}t\otimes u(t^{-1}v).
	\]
	反向对应为 \(F\mapsto\pi_W\circ F\)，其中 \(\pi_W(a\otimes w)=p_H(a)w\)，\(p_H:k[G]\to k[H]\) 保留群代数元素在 \(H\) 上的系数。这些对应均不要求 \([G:H]\) 在 \(k\) 中可逆。
\end{theorem}
\begin{proof}
	定义 \(p_H:k[G]\rightarrow k[H]\)，保留群代数元素在 \(H\) 上的系数，删去其余项。对 \(h_1,h_2\in H\)，有 \(p_H(h_1ah_2)=h_1p_H(a)h_2\)，因为左右乘 \(H\) 中元素均保持 \(H\) 及其补集。一个诱导分量 \(t\otimes W\) 应对应一个右陪集上的函数：将左陪集 \(tH\) 取逆得到 \(Ht^{-1}\)，并要求 \(t\otimes w\) 对应的函数在 \(ht^{-1}\) 处取 \(hw\)。条件 \(x\in Ht^{-1}\) 等价于 \(xt\in H\)，而 \(x=ht^{-1}\) 时恰有 \(xt=h\)。因此可以用 \(p_H(xt)\) 同时判断支集并给出作用于 \(w\) 的元素。对一般 \(a\in k[G]\) 线性延伸，令
	\[
		\Theta(a\otimes w)(x)=p_H(xa)w\qquad(x\in G).
	\]
	在 \(a=t\in G\) 时，这个公式具体就是
	\[
		\Theta(t\otimes w)(x)=
		\begin{cases}hw,&x=ht^{-1},\ h\in H,\\0,&x\notin Ht^{-1}.\end{cases}
	\]
	\(\Theta(a\otimes w)(hx)=hp_H(xa)w\)，故其值确在余诱导空间中。又 \(p_H(xah)w=p_H(xa)hw\)，所以它保持张量关系。\(G\) 等变性则来自
	\[
		\Theta(ga\otimes w)(x)=p_H(xga)w =\bigl(g\Theta(a\otimes w)\bigr)(x).
	\]

	选择左陪集代表 \(T\)。其逆元组成右陪集代表集 \(T^{-1}\)。对 \(s,t\in T\)，\(s^{-1}t\in H\) 当且仅当 \(s=t\)，所以
	\[
		\Theta(t\otimes w)(s^{-1})=
		\begin{cases}w,&s=t,\\0,&s\ne t.\end{cases}
	\]
	一个满足\eqref{b5:eq:coinduced-functions}的函数由它在 \(T^{-1}\) 上的值唯一确定。因此 \(\Theta\) 的逆为
	\[
		f\longmapsto\sum_{t\in T}t\otimes f(t^{-1}).
	\]
	这条和式因 \([G:H]<\infty\) 而有限；它在 \(T^{-1}\) 上逐个恢复函数值，因而恢复整个函数。\(\Theta\) 的定义本身不依赖 \(T\)，故逆映射也不依赖所选代表。对 \(H\)-同态 \(W\rightarrow W'\)，保留系数再作用与先送到 \(W'\) 的结果相同，证明自然性。

	再将\cref{b5:prop:coinduction-adjunction}中的 \(v\mapsto(x\mapsto u(xv))\) 与 \(\Theta^{-1}\) 复合，便得到所述 \(F_u\)。反向在 \(x=1\) 处求值，即为 \(\pi_W\circ F\)；\(p_H(ah)=p_H(a)h\) 保证 \(\pi_W\) 良定义，而 \(p_H(ha)=hp_H(a)\) 保证它是 \(H\)-线性的。这两条映射互逆且自然，已由所复合的两个对应证明。公式本身也显示代表选择无关：将 \(t\) 换为 \(th\)，对应项变成
	\[
		th\otimes u(h^{-1}t^{-1}v) =t\otimes h u(h^{-1}t^{-1}v) =t\otimes u(t^{-1}v).
	\]
	若 \(F(v)=\sum_{t\in T}t\otimes w_t(v)\)，且 \(t_0\in T\cap H\) 是单位陪集的代表，则 \(\pi_WF(v)=t_0w_{t_0}(v)\)；只有选 \(t_0=1\) 时才可直接说取出单位陪集的系数。整个构造没有任何除法。
\end{proof}

因此在有限指数情形，诱导既是限制的左伴随，也是其右伴随。左侧互反律把 \(W\to\operatorname{Res}_H^G V\) 延拓为 \(g\otimes w\mapsto gu(w)\)，右侧互反律把 \(\operatorname{Res}_H^G V\to W\) 延拓为上述陪集和式。右侧的求和使用陪集代表集的有限性，因而这个双重伴随结论适用于一般有限指数子群。

\ProseParagraphBreak
Etingof 等人的 Definition 5.8.1 直接使用\eqref{b5:eq:coinduced-functions}作为诱导模型，其 Theorem 5.10.1 因而首先给出右侧伴随公式\cite[Sections 5.8, 5.10, pp. 105--109]{EtingofEtAl2011RepresentationTheory}。\cref{b5:thm:induction-coinduction}说明了两种叙述的对应。对一般群，选左陪集代表 \(T\) 后，诱导的底层向量空间是 \(\bigoplus_{t\in T}W\)，而余诱导在右陪集代表 \(T^{-1}\) 上可任意取值，其底层向量空间是 \(\prod_{t\in T}W\)。有限指数时直和与直积相同；无限指数时，张量诱导对应的函数只在有限多个右 \(H\)-陪集上非零，余诱导却没有这一限制。这里的有限支集是相对于陪集空间 \(H\backslash G\) 而言；若 \(H\) 无限，函数在一个非零陪集内也可有无限多个非零值。逆映射的和式因而只对具有这种陪集有限支集的函数有定义。
